\section{Model}\label{sec:model}

The structural analysis of procurement preferences since \citet{marion2007} has modeled the set-aside as \emph{partial} preference: small and non-small firms both remain active in the auction, with one type awarded a bidding advantage---a price preference in \citeauthor{marion2007}, a scoring advantage with auction-level unobserved heterogeneity in \citet{krasnokutskaya2011}, an entry subsidy in \citet{atheyseira2011}. Under partial preference, the intensive margin of bidding and the extensive margin of entry are entangled by construction: the non-preferred type responds to both margins simultaneously, and the identification of one margin requires parametric assumptions about the other. The Brazilian reform studied in this paper is a qualitatively different regime. Non-SMEs leave the auction entirely; the SME pool re-equilibrates under a cleanly smaller bidder set; and the two margins become separately identified from drop-out data and observed post-policy participation counts, without parametric structure on the within-auction response. This is the analytical payload of the \emph{total-exclusion} design, and it is what makes a clean decomposition of sheltered bidding versus induced entry available here where it has not been available in partial-preference work.

The model formalizes the total-exclusion regime. $N$ potential bidders compete for a single indivisible contract. Each bidder $i$ has type $k_i \in \{\text{SME}, \neg\}$ and draws a private cost of supply $c_i$ from a type-specific distribution $F_c^k$, with $F_c^{\text{SME}}$ and $F_c^{\neg}$ differing in ways the data are allowed to determine.

The auction format is the identifying payload. BEC implements procurement through a Preg\~ao---a descending-clock electronic English-reverse auction, distinct both from first-price sealed-bid designs (Convite, in the smaller-value Brazilian segment) and from the Dutch descending-price format more commonly studied in the auction literature. The clock opens at the buyer's reference price $p^{\mathrm{ref}}$, firms post downward-revisable offers until only one firm remains, and that firm supplies the contract at the clearing price. What makes this format analytically convenient is a standard result of \citet{haile2003}: under independent private values, the weakly dominant strategy is to exit the auction when the running clock price reaches one's own cost. Drop-out prices of losers therefore \emph{point-identify} their costs, without requiring the researcher to invert a bid function or assume a parametric form for the bidding strategy. The winner, who never drops out before the clock stops, leaves only an upper bound $c_{\text{win}} \le c_{(2)}$ (the cost of the second-to-last firm to exit)---and that upper bound is exactly what the Turnbull robustness check in Section~\ref{app:robust} uses to deliver a non-parametric MLE of $F_c^k$. This point-identification under Preg\~ao is the reason a structural analysis is available here where it has not been available in cross-sectional work on auction preferences.

The model relies on three substantive restrictions, stated formally and defended empirically in the data.

\begin{assumption}[Independent Private Values, asymmetric across types]\label{a:ipv}
Costs $c_i$ are drawn independently across bidders. Conditional on type $k$, $c_i \sim F_c^k$ with density $f_c^k$ on support $[\underline{c}, \overline{c}] \subset \mathbb{R}_+$. The SME and non-SME distributions are not constrained to coincide.
\end{assumption}

Assumption~\ref{a:ipv} permits SMEs and non-SMEs to differ in cost distributions---a realistic feature given that non-SMEs are typically the large distributors of the same generic pharmaceutical---without requiring one to stochastically dominate the other.

\begin{assumption}[Multiplicative auction-level heterogeneity]\label{a:uh}
The normalized log-bid decomposes as $y_{it} = \log(b_{it} / p_t^{\mathrm{ref}}) = a_t + e_{it}$, where $a_t$ is an auction-level scale common to all bidders in auction $t$, and $e_{it}$ is bidder-specific and independent across bidders within $t$. Both $a_t$ and $e_{it}$ are mean-zero with finite variances $\sigma_a^{2,k}$ and $\sigma_e^{2,k}$ that may vary by stratum.
\end{assumption}

Assumption~\ref{a:uh} is the \citet{krasnokutskaya2011} restriction that separates a common scale from the bidder-specific private value. It is weaker than the fully non-parametric \citet{kotlarski1967} deconvolution, which would additionally require type-symmetry within $e_{it}$; Section~\ref{app:robust} reports a Turnbull NPMLE that dispenses with Assumption~\ref{a:uh} altogether and corroborates the baseline numbers within five percent.

\begin{assumption}[Set-aside, with equilibrium SME selection]\label{a:setaside}
Under the open regime both SMEs and non-SMEs enter with strictly positive probability. Under the SME-only regime, only SMEs are admissible. Within a short window around the cutoff, the underlying technology of supply is stable, but the equilibrium distribution of active SMEs under the restricted regime may differ from its pre-regime counterpart because the policy changes which SMEs find entry profitable. Write these distributions $F_c^{\mathrm{SME},\tau}$ for $\tau \in \{\mathrm{Pre},\mathrm{Post}\}$.
\end{assumption}

Assumption~\ref{a:setaside} is the structural counterpart of the parallel-trends assumption in reduced-form DiD, but with equilibrium selection made explicit. In the \citet{atheyseira2011} reading adopted here, the policy changes the admissible set and therefore the composition of the SME bidders who actually enter. The post-policy SME cost distribution is thus an equilibrium object, not evidence of a primitive technological break. This interpretation fits the diagnostic evidence: the primitive-invariance KS test in Table~\ref{tab:v3_primitive_invariance} rejects strict invariance precisely where one would expect entry selection to bite most strongly (pharma Preg\~ao SMEs), while it fails to reject in the strata where comparability is plausible (pharma Convite non-SMEs, $D = 0.032$; non-pharma non-SMEs, $D = 0.051$). The cross-modality check between Convite GPV and Preg\~ao drop-out recoveries in Table~\ref{tab:v3_cross_modality} provides an additional consistency test.

\paragraph{Strict invariance as a limiting case and as a falsification test.} A stronger reading would impose $F_c^{\text{SME,Post}} = F_c^{\text{SME,Pre}}$ and attribute the entire policy response to bidder counts alone. I treat that as the model's own falsification test of Assumption~\ref{a:setaside}, not as the main specification. The test is reported explicitly in Section~\ref{sec:falsification} alongside the headline welfare-dominance claim, and in Section~\ref{app:robust} as a structural-decomposition robustness. Three results follow. The intensive-margin share is preserved across specifications (75--85 percent under both). The non-pharma welfare-dominance of the 10 percent price preference is preserved across specifications. The pharma welfare comparison reverses: under strict invariance the implicit welfare weight required to prefer the set-aside falls to 0.7, below the utilitarian benchmark. The bifurcation is therefore confined to one stratum, identifiable in advance by goods characteristics (thin protected pool, high heterogeneity), and is itself substantive evidence on the role of equilibrium selection in pharma price formation.

\paragraph{Equilibrium and counterfactuals.} In the asymmetric IPV English-reverse auction, the winning price $p$ equals in expectation the second-lowest cost drawn from the active pool, $p = E[c_{(2)}]$. This is a direct consequence of the \citet{vickrey1961} revenue-equivalence argument as carried over to descending-clock settings and exploited by \citet{krasnokutskaya2011} and \citet{maskinriley2000}. The counterfactuals therefore reduce to Monte Carlo simulation over draws from the estimated $F_c^k$: draw a stratum-specific count of bidders from each type's Poisson arrival process, draw costs, sort, and read off the two lowest order statistics. The lowest, $c_{(1)}$, is the production cost; the second-lowest, $c_{(2)}$, is the expected transaction price.

\paragraph{Entry.} The BNE takes equilibrium arrival rates as primitives rather than solving for them explicitly. Observed Pre-period counts $(N^{\text{SME}}_{\text{Pre}}, N^{\neg}_{\text{Pre}})$ are the status-quo arrival rates under the open regime; Post-period counts $(N^{\text{SME}}_{\text{Post}}, 0)$ are the equilibrium arrival under the restricted regime. This is a reduced-form treatment of entry in the \citet{atheyseira2011} spirit---I do not estimate a structural fixed cost of entry jointly with the cost distributions---but it is sufficient for the counterfactual arithmetic. The main specification therefore combines observed Post counts with the observed Post SME bidder distribution. The implied zero-profit entry costs per bidder are reported separately (Table~\ref{tab:v3_entry_cost}) and are in the R\$0.55--R\$2.46 range, plausible given the overhead of preparing a BEC tender response.

\paragraph{The decomposition and \emph{sheltered bidding}.} Three counterfactual scenarios are of interest. $S_1$ fixes the regime as open and the pool as pre-period; it is the baseline. $S_2$ imposes SME-only while holding the SME pool at its pre-period size; this isolates the intensive effect of restricting the bidding set while shutting off the entry response. $S_3$ imposes SME-only at the observed post-period SME pool size; this is the full policy equilibrium. The total price effect $p_{S_3} - p_{S_1}$ decomposes cleanly:
\begin{equation}
\underbrace{p_{S_3} - p_{S_1}}_{\text{total effect}}
\;=\; \underbrace{(p_{S_2} - p_{S_1})}_{\text{intensive margin (sheltered bidding)}}
   \,+\, \underbrace{(p_{S_3} - p_{S_2})}_{\text{entry margin}}.
\end{equation}
Both components are estimated directly; neither is treated as a residual. The intensive margin---the \emph{sheltered bidding} effect of the introduction---captures the sharpening of the second-order statistic when non-SMEs are removed from a pool in which they were, on average, the more numerous type; the SME bidders who survive the exclusion respond to their sheltered position by bidding closer to the reference price. The entry margin captures the partially offsetting effect of SMEs doubling their participation rate in response to the policy. The two are the objects the partial-preference literature \citep{marion2007, krasnokutskaya2011, atheyseira2011} has struggled to separate because, under partial preference, both types remain in the auction and both margins enter each type's equilibrium response simultaneously. Under the total-exclusion design of the Brazilian reform, the non-SME pool is counterfactually zero in $S_2$ and $S_3$, and the decomposition above follows without additional parametric structure on the SME response. They are the central quantitative contribution of Section~\ref{sec:results}.
